\documentclass{article}
\usepackage{graphicx}
\usepackage{amsmath}
\usepackage{url}
\usepackage{hyperref}
\usepackage[margin=0.8in]{geometry}
\usepackage{listings}
\usepackage{enumitem}
\setlist{itemsep=2pt,topsep=4pt}
\lstset{basicstyle=\small\ttfamily,breaklines=true,columns=fullflexible,
  keepspaces=true,showstringspaces=false,aboveskip=4pt,belowskip=4pt}

\title{Project Milestone 1: Verifiable Tasks}
\author{COMP 480/580, Fall 2026}
\date{Due: October 31}

\begin{document}
\maketitle

\section{Introduction}

AI agents use tools to retrieve information, process data, and create files.
An agent's claim that it has finished a task, however, does not show that its
result is correct. To evaluate an agent reliably, we need a precise task
specification, a reproducible environment, and an executable verifier that
checks the actual result. In this milestone, you will create an original task,
its execution environment, and a verifier that checks its outcome. Use the
public starter below to build and validate your task locally, then publish it
on \href{https://rush-lab.org}{RUSH Lab}.

\section{Tasks}

\begin{enumerate}
  \item \textbf{Task:} Define its inputs, tools, outputs, and success
    conditions. Renaming the starter is not enough.
  \item \textbf{Environment:} Include the data, dependencies, and reproducible
    setup instructions.
  \item \textbf{Verifier:} Check actual results. Test correct, incorrect,
    missing, malformed, and boundary cases. \textbf{The verifier must not call
    an LLM.} You may use an LLM to help write it, but the verifier itself is a
    program that runs without one.
  \item \textbf{Validation and publication:} Test, run, and publish the task,
    and provide reproducible evidence.
\end{enumerate}

\section{Instructions}

\subsection{Set up your plugin}

Install Git, Python 3.12 or later, and \texttt{uv}. Start from the
\href{https://github.com/zeyuyang8/arxiv-source}{arXiv Search and Source
Downloader plugin}, the reference plugin: clone it under a new name, with a
history of its own, and check that its tests pass:

\begin{lstlisting}
git clone https://github.com/zeyuyang8/arxiv-source.git my-task
cd my-task
rm -rf .git && git init
uv sync --extra dev
uv run pytest
\end{lstlisting}

\noindent
Replace \texttt{my-task} with a unique id in lowercase words joined by
hyphens. Set it as the \texttt{id} in \texttt{nanoagent.yaml}, with a
\texttt{name} of your own, and name the project \texttt{nanoagent-my-task} in
\texttt{pyproject.toml}: RUSH Lab refuses an id another plugin already has.
Its README and source show how the task, tools, environment, and verifier fit
together. Replace them with your own original task: its manifest, contract,
tools, and verifier. Keep the required package structure, and keep reports and
run outputs outside the plugin folder.

\medskip
\noindent
In an agent task, tools expose the inputs and save the agent's answer to a
workspace file; the agent does the reasoning. An optional \texttt{reference}
tool computes the answer key before the agent runs, and it must not be one of
the agent's task tools. The verifier receives this pinned reference. If you
have no reference tool, set \texttt{reference: null} and explain how
correctness is established.

\subsection{Implement and validate}

The verifier implements \texttt{verify\_input}, \texttt{verify\_output}, and
\texttt{verify\_task}. The first two check tool calls; the third checks the
completed task. Each returns \texttt{valid}, a \texttt{score} in $[0,1]$, a
\texttt{message}, and \texttt{details}. Follow the starter's signatures and
typed inputs. The verifier must be deterministic: no LLM calls, network
access, randomness, subprocesses, or file writes.

\medskip
\noindent
Keep the dependencies and \texttt{uv.lock} current, keep the generated
setup, UI, and product files, and list only variable names in \texttt{.env}.
Set \texttt{status: ready} in \texttt{nanoagent.yaml}. Then, from your plugin
folder, run:

\begin{lstlisting}
uv sync --extra dev
uv run pytest
\end{lstlisting}

\noindent
Fix all errors. Put your verifier tests in \texttt{tests/test\_verifier.py}.

\subsection{Publish and validate on RUSH Lab}

Send your GitHub username to the course staff for an invitation, then sign in
at \url{https://rush-lab.org} with the invited GitHub account.

\medskip
\noindent
The agent and your plugin run on your own computer, in RUSH Lab's background
service; you keep using RUSH Lab in your browser. To install it, open
\textbf{Settings $\rightarrow$ AI}, copy the command shown there, and paste it
in a terminal. It connects your computer to your account by itself. It runs on
macOS (Apple silicon or Intel) and on Linux (x86-64), where it also needs
\texttt{bubblewrap} (\texttt{sudo apt install bubblewrap} on Ubuntu). On
Windows, install WSL with \texttt{wsl --install} in PowerShell, open Ubuntu,
install \texttt{bubblewrap} there, and paste the command in Ubuntu; keep that
window open while you use RUSH Lab. The command is yours alone: do not share it.
If it has no OpenRouter API key, RUSH Lab gives it one;
you can instead save your own key from
\href{https://openrouter.ai/settings/keys}{OpenRouter} under the same pane.
RUSH Lab sets the model.

\medskip
\noindent
Push the complete, tested plugin to a \textbf{public GitHub repository}. Open
the \href{https://rush-lab.org/w/comp-480-580/marketplace}{Marketplace},
select \textbf{Submit Your Artifact}, paste the repository URL, tick the
statement, and select \textbf{Submit}. If the plugin is in a subfolder, use
its \path{/tree/<branch>/<folder>} URL. Fix any reported errors and submit
again.

\medskip
\noindent
\textbf{Passing the publication checks confirms the package format only; it
does not show that your task or verifier is correct.}

\medskip
\noindent
Ask a workspace owner to add your plugin, then open it, ask its chat to run
your task, and save the task verifier's verdict shown with the result. Have
another student or a TA reproduce your setup and your correct- and
incorrect-result checks from the README, and record any issues and fixes.

\section{Questions}

In your report, explain:
\begin{enumerate}
  \item why your task is objectively verifiable;
  \item how your environment is reproducible;
  \item how the verifier handles incorrect outputs and boundary cases; and
  \item what limitations remain.
\end{enumerate}

\section{What to Submit}

Submit the GitHub URL and commit hash; your RUSH Lab plugin
identifier; the complete source and lockfile; the tests; and a README with
exact setup and execution instructions. Outside the plugin folder, include the
inputs, validation logs, verdicts, and the independent reproduction record.
Also submit a short report answering the questions above. TAs must be able to
reproduce your results without manual fixes.

\end{document}
